% Shared preamble of main_zh.tex and supplement_zh.tex.
\usepackage[margin=1in]{geometry}
\usepackage{amsmath,amssymb}
\usepackage{graphicx}
\graphicspath{{../}}
\usepackage{booktabs}
\usepackage{tabularx}
\usepackage{multirow}
\usepackage{longtable}
\usepackage{array}
\usepackage{xcolor}
\usepackage[most]{tcolorbox}
\usepackage{enumitem}
\usepackage[font=small,labelfont=bf]{caption}
\usepackage{subcaption}
\usepackage[numbers,sort&compress]{natbib}
\usepackage{xurl}
\usepackage[section]{placeins}
\usepackage[bookmarks=true,colorlinks=true,linkcolor=blue!50!black,citecolor=blue!50!black,urlcolor=blue!50!black]{hyperref}
\hypersetup{pdftitle={类型化决策模型能否承担 PoL 治理中的自动化保障职能？——基于早期证据的爱2证明（PoL2）规范逐条分析},pdfauthor={鄢申涛}}
\usepackage{cleveref}
% cleveref in Chinese: one helper sets every format for a label type.
\newcommand{\zhcref}[3]{%
  \crefformat{#1}{##2#2~##1~#3##3}\Crefformat{#1}{##2#2~##1~#3##3}%
  \crefrangeformat{#1}{##3#2~##1~#3##4至##5#2~##2~#3##6}\Crefrangeformat{#1}{##3#2~##1~#3##4至##5#2~##2~#3##6}%
  \crefmultiformat{#1}{##2#2~##1~#3##3}{和##2#2~##1~#3##3}{、##2#2~##1~#3##3}{和##2#2~##1~#3##3}%
  \Crefmultiformat{#1}{##2#2~##1~#3##3}{和##2#2~##1~#3##3}{、##2#2~##1~#3##3}{和##2#2~##1~#3##3}%
  \crefrangemultiformat{#1}{##3#2~##1~#3##4至##5#2~##2~#3##6}{和##3#2~##1~#3##4至##5#2~##2~#3##6}{、##3#2~##1~#3##4至##5#2~##2~#3##6}{和##3#2~##1~#3##4至##5#2~##2~#3##6}%
  \Crefrangemultiformat{#1}{##3#2~##1~#3##4至##5#2~##2~#3##6}{和##3#2~##1~#3##4至##5#2~##2~#3##6}{、##3#2~##1~#3##4至##5#2~##2~#3##6}{和##3#2~##1~#3##4至##5#2~##2~#3##6}}
\zhcref{section}{第}{节}
\zhcref{subsection}{第}{节}
\zhcref{subsubsection}{第}{节}
\zhcref{appendix}{附录}{}
\zhcref{figure}{图}{}
\zhcref{table}{表}{}
\zhcref{enumi}{第}{项}
\zhcref{supplement}{补充材料}{}
\definecolor{f4pblue}{HTML}{0F4D92}
\definecolor{f4pred}{HTML}{B64342}
\definecolor{f4pgreen}{HTML}{8BCF8B}
% CJK passages need no special handling under xelatex/ctex.
\newcommand{\zh}[1]{#1}
\newcommand{\axis}[1]{\textbf{#1}}
\newcommand{\jev}{\textsc{Jev}}
\newcommand{\pol}{PoL2}
\newcommand{\clause}[1]{\S\,#1}
\newcommand{\Clause}[1]{\S\,#1}
\newcommand{\art}[1]{第~#1~条}
\newcommand{\Art}[1]{第~#1~条}
% 来源标记：Zenodo 记录（未经预印本服务器审核）；灰色文献（博客、厂商页面、代码仓库）
\newcommand{\zmark}{\textsuperscript{$\ddagger$}}
\newcommand{\gmark}{\textsuperscript{$\dagger$}}
% Generated by scripts/merge_coding.py; do not edit by hand.
\newcommand{\nArxiv}{79}
\newcommand{\nZenodo}{27}
\newcommand{\nOther}{28}
\newcommand{\nPapers}{107}
\newcommand{\nGrey}{7}
\newcommand{\nCoded}{85}
\newcommand{\nRows}{161}
\newcommand{\nS}{16}
\newcommand{\nQ}{97}
\newcommand{\nN}{48}
\newcommand{\nPolClauses}{28}
\newcommand{\nKappaN}{30}
\newcommand{\nKappaAgree}{26 of 30}
\newcommand{\nKappa}{0.72}
\newcommand{\nKappaLo}{0.42}
\newcommand{\nKappaHi}{0.94}
\newcommand{\nKappaMarg}{S 3, Q 23, N 4}
\newcommand{\nStrongS}{7}
\newcommand{\nStrongQ}{45}
\newcommand{\nStrongN}{27}
\newcommand{\nStrongRows}{79}
\newcommand{\nNoZenS}{13}
\newcommand{\nNoZenQ}{75}
\newcommand{\nNoZenN}{29}
\newcommand{\nNoZenRows}{117}
\newcommand{\nStronger}{37}
\newcommand{\nModerate}{17}
\newcommand{\nWeaker}{38}
\newcommand{\nAppraised}{92}

% Two counts are English phrases in counts.tex; Chinese forms here (keep in sync).
\renewcommand{\nKappaAgree}{30 对中的 26 对}
\renewcommand{\nKappaMarg}{S 3、Q 23、N 4}
% Generated by scripts/merge_update.py (2026-10-08 update search); do not edit by hand.
\newcommand{\nUDocs}{57}
\newcommand{\nUCoded}{48}
\newcommand{\nUArxiv}{27}
\newcommand{\nUZenodo}{21}
\newcommand{\nURows}{118}
\newcommand{\nUS}{12}
\newcommand{\nUQ}{92}
\newcommand{\nUN}{14}
\newcommand{\nUStrongRows}{59}
\newcommand{\nUStrongS}{5}
\newcommand{\nUStrongQ}{45}
\newcommand{\nUStrongN}{9}
\newcommand{\nUKappaN}{118}
\newcommand{\nUKappaAgree}{97 of 118}
\newcommand{\nUKappa}{0.55}
\newcommand{\nUKappaLo}{0.36}
\newcommand{\nUKappaHi}{0.70}
\newcommand{\nUSetPostWindow}{30}
\newcommand{\nUSetLateIndexed}{5}
\newcommand{\nUSetReScreened}{13}

% English phrase in counts_update.tex; Chinese form here (keep in sync).
\renewcommand{\nUKappaAgree}{118 对中的 97 对}
\newcolumntype{L}{>{\raggedright\arraybackslash}X}
\setlist{itemsep=2pt,topsep=4pt}
\newtcolorbox{answerbox}{breakable,enhanced,colback=f4pblue!5,colframe=f4pblue!60,boxrule=0.6pt,
  arc=1.5pt,left=6pt,right=6pt,top=4pt,bottom=4pt,before skip=6pt,after skip=8pt,
  title={\small\bfseries 简要回答},fonttitle=\small\bfseries,coltitle=white,colbacktitle=f4pblue!80}
\renewcommand{\abstractname}{摘\quad 要}
